\documentclass[12pt]{article}
\usepackage{notes-project}
\usepackage{graphicx}
\usepackage{subcaption}
\usepackage{listings}

\graphicspath{{lec03/}}
\numberwithin{figure}{section}

\lstset{
  language=Python,
  basicstyle=\ttfamily\footnotesize,
  frame=single,
  framesep=6pt,
  xleftmargin=6pt,
  xrightmargin=6pt,
  breaklines=true,
  breakatwhitespace=true,
  showstringspaces=false,
  columns=fullflexible,
  keepspaces=true,
  keywordstyle=\color{blue!60!black},
  commentstyle=\color{gray},
  stringstyle=\color{green!40!black},
}

\begin{document}

\begin{center}
  {\LARGE\bfseries Lecture 3: Numerical Differentiation}\\[0.15em]
  {\large Introduction to Data Analysis Techniques}\\[0.15em]
\end{center}

\section{Taylor Approximation}

\subsection{Taylor Polynomial and Remainder}

\missingprereq{Continuity, differentiability, and smooth function spaces,
real-analysis}

We write $\SmoothRealFunctions$ for functions $\RR\to\RR$ whose derivatives
of every order, including order zero, are continuous

\begin{dfn}{taylor-polynomial-generator}
  The Taylor term constructor takes $f$, then a centre $a$, then an order
  $n$, and returns the $n$th term as a polynomial function
  \[
    \begin{aligned}
      \mathrm{TaylorGen} &: \SmoothRealFunctions\to\RR\to\NN
        \to(\RR\to\RR)\\
      &\defeq f\map\bigl(a\map\bigl(n\map\bigl(x\map
        \frac{\DiffOp^n(f)(a)}{n!}\cdot(x-a)^n\bigr)\bigr)\bigr)
    \end{aligned}
  \]
\end{dfn}

\begin{dfn}{taylor-polynomial}
The $n$-th order Taylor polynomial generator is defined as 
\[\mathrm{TaylorPoly} : \SmoothRealFunctions\to\RR\to\NN\to\RR\to\RR \defeq f \map a \map N \map x \map \sum^N n \map \mathrm{TaylorGen}(f)(a)(n)(x) \]
\end{dfn}

\begin{conv}{taylor-polynomial-notation}

\end{conv}

\begin{dfn}{taylor-remainder}

\end{dfn}

\begin{thm}[Finite Taylor Expansion]{finite-taylor-expansion}
  Let $f\in\SmoothRealFunctions$, $n\in\NN$, and $a\in I$, where
  $I\subseteq\RR$ is an open interval. For the
  \xref[Taylor polynomial]{def:taylor-polynomial},
  \[
    \begin{gathered}
      \forall x\in I\setminus\{a\},\quad
      \exists\xi\in\bigl(\min(a,x),\max(a,x)\bigr):\\
      f(x)=P_{f,a,n}(x)+\frac{\DiffOp^{n+1}(f)(\xi)}{(n+1)!}
        \cdot(x-a)^{n+1}
    \end{gathered}
  \]
  At $x=a$, $f(a)=P_{f,a,n}(a)$ and the remainder is zero
\end{thm}

\subsection{Taylor Series}

% Moved from lec01.tex. Preserve these semantic IDs for references to the
% Taylor material in this document.
\begin{dfn}[Taylor Series]{taylor-series}
  Let $f\in\SmoothRealFunctions$ and $a\in\RR$.
  Its \emph{Taylor-series sum} at $x$, when it exists, is
  \[
    f(a)+\infsum\bigl(n\map\TaylorGen(f)(a)(n)(x)\bigr)
  \]
  The sequence is indexed by $\NNOne$. The displayed expression equals
  $f(x)$ precisely when $R_{f,a,N}(x)\to0$ as $N\to\infty$
\end{dfn}

\begin{rem}[Taylor Series Approximation]{taylor-series-approx}
  We can truncate the \xref[Taylor series]{def:taylor-series} to obtain a
  finite approximation. The
  \xref[Taylor expansion with remainder]{thm:finite-taylor-expansion}
  expresses its error under suitable differentiability assumptions

  Adding terms does not automatically improve the approximation. The
  remainder must be controlled; equality with the infinite series requires
  that the remainder tend to zero as the number of terms increases
\end{rem}

\subsection{Example: Electric Field of a Dipole}

\begin{eg}[Electric Field of a Dipole for $x \gg a$]{dipole-field}
  Consider the electric field on the axis of a dipole with charges $+q$ at
  $x=a$ and $-q$ at $x=-a$, where $a>0$ and $x>a$
  \[
    E=\frac{q}{4\pi\varepsilon_0}\cdot\left[
      \frac{1}{(x-a)^2}-\frac{1}{(x+a)^2}
    \right]
  \]

  Expanding in $a/x$ using the \xref[Taylor series]{def:taylor-series} gives
  \[
    E=\frac{q}{4\pi\varepsilon_0\cdot x^2}\cdot\left[
      4\frac{a}{x}+8\left(\frac{a}{x}\right)^3
      +12\left(\frac{a}{x}\right)^5+\cdots
    \right]
  \]

  Keeping the first two nonzero terms gives an approximation for $x\gg a$
  \[
    E\approx\frac{q}{4\pi\varepsilon_0\cdot x^2}\cdot\left[
      4\frac{a}{x}+8\left(\frac{a}{x}\right)^3
    \right]
  \]
  The leading term is proportional to $x^{-3}$
\end{eg}

\subsection{Taylor Expansion in Two Variables}

\todo[Write the expansion through second order, including the mixed derivative]

\section{Error Orders}

\todo[Define big-O notation as the step size tends to zero; explain how halving the step changes an error of order one or two]

\section{Numerical Differentiation}

\subsection{Forward and Backward Differences}

\todo[Define both first-derivative approximations and derive their error orders from Taylor expansion]

\subsection{Central Differences}

\todo[Derive central approximations for the first and second derivatives and their error orders]

\subsection{Programming Exercise}

\todo[Implement the difference formulas for sine at x=1; compare with the exact derivatives over several step sizes]

\section{First-Order Differential Equations}

\todo[Define an ordinary differential equation and an initial-value problem in separate blocks]

\section{Euler's Method}

\todo[Introduce the time grid and numerical iterates; derive the forward Euler update and distinguish local and global errors]

\subsection{Programming Exercise}

\todo[Solve x'=-x with x(0)=1 and compare with exp(-t) using several time steps]

\section{Second-Order Dynamics}

\todo[Introduce velocity to rewrite a second-order equation as two first-order equations; derive Euler updates using the old values on both right-hand sides]

\subsection{Programming Exercise}

\todo[Implement the pendulum model and compare results as the time step changes]

\end{document}
